\documentclass{situationist-preprint}

\title{Bracketing the Bracket}
\msid{PHE-01-b}
\msfamily{PHE-01}
\msdiscipline{Phenomenology}
\mstype{Companion manuscript}
\mskeywords{epoché; natural attitude; reflection; iteration; transcendental ego}

\begin{document}
\maketitle

\begin{abstract}
PHE-01-a argues that the epoché can be sustained as a reflective attitude over finite intervals but never permanently, because the reflective attitude lasts only as long as acts of bracketing are renewed. A natural proposal for escaping this result is to bracket the bracket: to perform a higher-order epoché that suspends the act of suspension itself, either to make the reduction self-sustaining or to make it complete. We argue that the proposal fails in both forms. A bracketing of the bracket is a further act, distinct from and later than the act it brackets; it is subject to the same renewal and lapse; and it leaves its own act unbracketed. Iteration multiplies acts without extending the duration of the epoché.
\end{abstract}

\section{The proposal}

PHE-01-a derives its main result (Theorem~4.5 there) from three premises: the reflective attitude obtains at a moment exactly when an act of bracketing was performed at the preceding moment, while the natural attitude requires no act (Assumption~3.5); renewal of such acts lapses within a bounded span $L$ (Assumption~3.6); and within that span renewal is possible (Assumption~3.7). It follows that the epoché can be sustained over every interval of at most $L$ moments, is permanent in no course of conscious life, and that the natural attitude returns within $L+1$ moments of any moment.

The proposal we examine accepts the premises but denies that they apply to every kind of bracketing. It comes in two forms. On the \emph{stabilizing} form, a second-order epoché suspends the first-order act of bracketing, and with it the need to renew that act: once the bracketing is itself bracketed, the reflective attitude would no longer depend on anything that could lapse. Fink's account of the \enquote{transcendental onlooker}, who does not participate in the positings he observes, has been read as a picture of such a stabilized reduction \citep{Fink1933}. On the \emph{completing} form, the second-order epoché is needed because the first-order one is incomplete: it brackets everything except itself, and only by bracketing the bracket does the reduction become universal.

\section{The argument}

\begin{proposition}\label{p:distinct}
Every bracketing of an act of bracketing is an act distinct from, and later than, the act it brackets.
\end{proposition}

\begin{proof}
No act of bracketing brackets itself (Assumption~2.2 of PHE-20-a). Hence whatever brackets an act $b$ is an act other than $b$. It is later than $b$, since one cannot suspend the positing character of an act that has not yet been performed: what is bracketed is given in reflection, and reflection is directed at an act already under way \citep[\S77]{Husserl1913}.
\end{proof}

\begin{proposition}\label{p:nolonger}
Iterated bracketing does not extend the bound of Theorem~4.5 of PHE-01-a.
\end{proposition}

\begin{proof}
A bracketing of a bracketing either keeps the general thesis out of action or it does not. If it does not, it does not sustain the reflective attitude at all and is irrelevant to its duration. If it does, it is an act of bracketing in the sense of Definition~3.2 of PHE-01-a, directed in addition at the first-order act. By Proposition~\ref{p:distinct} it is a further act, and the reflective attitude it sustains obtains only at the moment after it is performed (Assumption~3.5 of PHE-01-a). Assumption~3.6 bounds runs of renewal whatever the order of the acts renewed, since it concerns the performance of acts, not their objects. A course of conscious life in which first-order and second-order acts alternate, or in which only second-order acts are performed, therefore contains a lapse within every span of $L+1$ moments, exactly as before.
\end{proof}

Proposition~\ref{p:nolonger} disposes of the stabilizing form. The second-order act does not relieve the subject of the need to renew; it is one more thing that needs renewing. If there were a bracketing that did not need to be renewed, it would not be an act but a standing condition, and by Assumption~3.5 of PHE-01-a the only attitude that obtains without any act is the natural attitude (Definition~3.1), in which the general thesis is in force. The transcendental onlooker is therefore not a stabilized state but a name for what is performed whenever the reduction is performed.

\begin{proposition}\label{p:incomplete}
No finite or infinite series of bracketings brackets every act in it.
\end{proposition}

\begin{proof}
If the series is finite, its last act is bracketed by no act of the series, by Proposition~\ref{p:distinct}. If it is infinite, each act is bracketed by a later one, but the series is never completed at any moment of a conscious life, and at each moment the latest act performed is unbracketed. Either way, something remains: the residue of Theorem~3.1 of PHE-20-a.
\end{proof}

Proposition~\ref{p:incomplete} disposes of the completing form. The incompleteness that motivates it is not a defect to be remedied by a higher order; it reappears at every order. This is the result of PHE-20-a, and it is also the point on which PHE-22-a turns: whatever brackets, brackets everything except the bracketing itself.

\section{Objections}

\begin{objection}
Proposition~\ref{p:nolonger} assumes that the lapse of Assumption~3.6 of PHE-01-a is indifferent to the order of the acts. But higher-order acts may be less effortful, as habits are, and might lapse more slowly.
\end{objection}

\begin{reply}
That would change the value of $L$, not the form of the result. A longer bound still yields an interval after which the natural attitude returns. Moreover, the higher-order act is not easier than the first-order one: it presupposes it, since it must find a first-order bracketing to suspend.
\end{reply}

\begin{objection}
Bracketing the bracket is a way of returning to the natural attitude. If so, the proposal does not extend the epoché but describes its end.
\end{objection}

\begin{reply}
It does not describe its end either. To bracket the bracket is to suspend the positing that one is reflecting, which is itself a reflective accomplishment. The natural attitude, by Assumption~3.5 of PHE-01-a, returns without any act, simply because renewal stops. The return is by default, not by a counter-epoché; Merleau-Ponty's remark that the reduction cannot be completed concerns this default, not a further operation \citep{MerleauPonty1945}.
\end{reply}

\section{Conclusion}

Bracketing the bracket is always possible, and it is never more than another bracketing. It neither stabilizes the epoché nor completes it. The duration of the reflective attitude remains bounded as PHE-01-a states, and the unbracketed residue remains as PHE-20-a states. The two results are two sides of one fact: the epoché exists only in the acts that perform it.

\begin{thebibliography}{99}\small
\bibitem[Fink(1933)]{Fink1933} Fink, E. (1933). Die phänomenologische Philosophie Edmund Husserls in der gegenwärtigen Kritik. \emph{Kant-Studien} 38: 319--383.
\bibitem[Husserl(1913)]{Husserl1913} Husserl, E. (1913). Ideen zu einer reinen Phänomenologie und phänomenologischen Philosophie. Erstes Buch. \emph{Jahrbuch für Philosophie und phänomenologische Forschung} 1.
\bibitem[Merleau-Ponty(1945)]{MerleauPonty1945} Merleau-Ponty, M. (1945). \emph{Phénoménologie de la perception}. Gallimard.
\end{thebibliography}

\end{document}
